% ------------------------------------------------------------------------
% LaTeX - Preambel ******************************************************
% ------------------------------------------------------------------------
% pre-work
% ========================================================================
% % ToDo kennzeichnen
\newcommand{\workTodo}[1]{\textcolor{red}{todo: #1}}

% % Für Datum und Zeit in Fusszeile
% % !!!Inhalt bei Fertigstellung der Arbeit löschen
\newcommand{\workMarkDateTime}{\workTodo{\today{} - \thistime\ Uhr}}

% % Alle Namen werden im Titel und im hyperref-Paket eingetragen
% % !!! Ueberall für <Wert> das Entsprechende eintragen

 % <Typ> Studienarbeit, Dipolmarbeit, Studienarbeit oder Bachlor-Abschlussarbeit
\newcommand{\workTyp}{\workTodo{<Typ>}\xspace}

 % <Titel> der Arbeit
\newcommand{\workTitel}{\workTodo{<Titel>}}

 % <Studiengang> z.B. Kommunikationstechnik
\newcommand{\workStudiengang}{\workTodo{<Studiengang>}\xspace}

% <Semester> mit Jahr z.B. Sommersemester 2008  
\newcommand{\workSemester}{\workTodo{<Semester>}\xspace}

% <Name> des Studenten
\newcommand{\workNameStudent}{\workTodo{<Name>}\xspace}

% <Pruefer> Name des prüfenden (betreuenden) Professor an der Hochschule
\newcommand{\workPruefer}{\workTodo{<Pruefer>}\xspace} 


% %%% Nur bei Abschluss-Arbeiten

% <Datum> der Abgabe der Arbeit (Eidesstatliche Erklärung)
\newcommand{\workDatum}{\today\xspace}

% <Zweitprüfer>
\newcommand{\workZweitPruefer}{\workTodo{<Zweitprüfer>}\xspace}

% <Zeitraum>
\newcommand{\workZeitraum}{\workTodo{<Zeitraum>}\xspace}


% %%% Nur bei Industrie-Arbeiten:

% <Firma>
\newcommand{\workFirma}{\workTodo{<Firma>}\xspace}

% <Betreuer in der Firma>
\newcommand{\workBetreuer}{\workTodo{<Betreuer in der Firma>}\xspace}

% Firmenlogo Name hier anpassen, Größe (wenn möglich) nicht ändern
\newcommand{\workFirmenLogo}{\includegraphics[width=5cm]{fig/aa-titel/GNU_logo.eps}} 
